\import{relation/transitive-closure.tex}

\section{Joinability and confluence}

\begin{definition}\label{joinability_relation}
    Let $R$ be a binary relation on $X$.
    $\joinabilityRelation{X}{R}=
        \converse{\reflexiveTransitiveClosure{X}{R}}
            \circ\reflexiveTransitiveClosure{X}{R}$.
\end{definition}

\begin{classification}\label{joinability_relation_type}
    Let $R$ be a binary relation on $X$.
    Then $\joinabilityRelation{X}{R}$ is a binary relation on $X$.
\end{classification}
\begin{proof}
    $\reflexiveTransitiveClosure{X}{R}$ is a binary relation on $X$
        by notion.
    $\converse{\reflexiveTransitiveClosure{X}{R}}$ is a binary relation on
        $X$ by notion.
    $\converse{\reflexiveTransitiveClosure{X}{R}}
        \circ\reflexiveTransitiveClosure{X}{R}$ is a binary relation on $X$
        by notion.
    Follows by definition.
\end{proof}

\begin{definition}\label{joinable_on}
    Let $R$ be a binary relation on $X$.
    Let $b,c\in X$.
    $b$ is joinable with $c$ in $X$ with respect to $R$ iff
        $b\mathrel{\joinabilityRelation{X}{R}}c$.
\end{definition}

\begin{proposition}\label{joinability_relation_iff}
    Let $R$ be a binary relation on $X$.
    Let $b,c\in X$.
    Then $b$ is joinable with $c$ in $X$ with respect to $R$ iff
        there exists $d\in X$ such that
            $b\mathrel{\reflexiveTransitiveClosure{X}{R}}d$ and
            $c\mathrel{\reflexiveTransitiveClosure{X}{R}}d$.
\end{proposition}
\begin{proof}
    Let $T=\reflexiveTransitiveClosure{X}{R}$.
    $T$ is a binary relation on $X$ by notion.
    $T\subseteq X\times X$ by definition.
    Show that if $b$ is joinable with $c$ in $X$ with respect to $R$, then
        there exists $d\in X$ such that $b\mathrel{T}d$ and $c\mathrel{T}d$.
    \begin{subproof}
        Suppose $b$ is joinable with $c$ in $X$ with respect to $R$.
        Then $b\mathrel{(\converse{T}\circ T)}c$
            by \auto{joinable_on,joinability_relation}.
        Take $d$ such that $b\mathrel{T}d$ and
            $d\mathrel{\converse{T}}c$ by \obv{circ_iff}.
        $c\mathrel{T}d$ by \obv{converse_iff}.
        $(b,d)\in X\times X$ by \obv{subseteq}.
        Hence $d\in X$ by \obv{times_tuple_elim}.
    \end{subproof}
    Show that if there exists $d\in X$ such that $b\mathrel{T}d$ and
        $c\mathrel{T}d$, then
        $b$ is joinable with $c$ in $X$ with respect to $R$.
    \begin{subproof}
        Suppose there exists $d\in X$ such that $b\mathrel{T}d$ and
            $c\mathrel{T}d$.
        Take $d\in X$ such that $b\mathrel{T}d$ and $c\mathrel{T}d$.
        $d\mathrel{\converse{T}}c$ by \obv{converse_iff}.
        Hence $b\mathrel{(\converse{T}\circ T)}c$ by \auto{circ_iff}.
        Follows by \obv{joinable_on,joinability_relation}.
    \end{subproof}
    Follows by \auto.
\end{proof}

\begin{definition}\label{confluent_at_on}
    Let $R$ be a binary relation on $X$.
    Let $a\in X$.
    $R$ is confluent at $a$ on $X$ iff for all $b,c\in X$ such that
        $a\mathrel{\reflexiveTransitiveClosure{X}{R}}b$ and
        $a\mathrel{\reflexiveTransitiveClosure{X}{R}}c$
        we have $b$ is joinable with $c$ in $X$ with respect to $R$.
\end{definition}

\begin{definition}\label{locally_confluent_at_on}
    Let $R$ be a binary relation on $X$.
    Let $a\in X$.
    $R$ is locally confluent at $a$ on $X$ iff for all $b,c\in X$ such that
        $a\mathrel{R}b$ and $a\mathrel{R}c$
        we have $b$ is joinable with $c$ in $X$ with respect to $R$.
\end{definition}

\begin{definition}\label{confluent_on}
    Let $R$ be a binary relation on $X$.
    $R$ is confluent on $X$ iff
        for every $a\in X$ we have $R$ is confluent at $a$ on $X$.
\end{definition}

\begin{definition}\label{locally_confluent_on}
    Let $R$ be a binary relation on $X$.
    $R$ is locally confluent on $X$ iff
        for every $a\in X$ we have $R$ is locally confluent at $a$ on $X$.
\end{definition}

\begin{proposition}\label{joinability_relation_is_reflexive}
    Let $R$ be a binary relation on $X$.
    Then $\joinabilityRelation{X}{R}$ is reflexive on $X$.
\end{proposition}
\begin{proof}
    $\reflexiveTransitiveClosure{X}{R}$ is reflexive on $X$
        by \obv{reflexive_transitive_closure_is_reflexive}.
    Show that for every $b\in X$ we have
        $b\mathrel{\joinabilityRelation{X}{R}}b$.
    \begin{subproof}
        Fix $b\in X$.
        $b\mathrel{\reflexiveTransitiveClosure{X}{R}}b$
            by \obv{reflexive_on}.
        $b\mathrel{\converse{\reflexiveTransitiveClosure{X}{R}}}b$
            by \obv{converse_iff}.
        $b\mathrel{(\converse{\reflexiveTransitiveClosure{X}{R}}\circ
            \reflexiveTransitiveClosure{X}{R})}b$ by \auto{circ_iff}.
        Follows by \obv{joinability_relation}.
    \end{subproof}
    Follows by \obv{reflexive_on}.
\end{proof}

\begin{proposition}\label{joinability_relation_is_symmetric}
    Let $R$ be a binary relation on $X$.
    Then $\joinabilityRelation{X}{R}$ is symmetric.
\end{proposition}
\begin{proof}
    $\joinabilityRelation{X}{R}$ is a binary relation on $X$ by notion.
    $\joinabilityRelation{X}{R}\subseteq X\times X$ by definition.
    Show that for all $b,c$ such that
        $b\mathrel{\joinabilityRelation{X}{R}}c$ we have
        $c\mathrel{\joinabilityRelation{X}{R}}b$.
    \begin{subproof}
        Fix $b,c$.
        Suppose $b\mathrel{\joinabilityRelation{X}{R}}c$.
        $(b,c)\in X\times X$ by \obv{subseteq}.
        Hence $b,c\in X$ by \obv{times_tuple_elim}.
        $b$ is joinable with $c$ in $X$ with respect to $R$
            by \obv{joinable_on}.
        Take $d\in X$ such that
            $b\mathrel{\reflexiveTransitiveClosure{X}{R}}d$ and
            $c\mathrel{\reflexiveTransitiveClosure{X}{R}}d$
            by \obv{joinability_relation_iff}.
        $d\mathrel{\converse{\reflexiveTransitiveClosure{X}{R}}}b$
            by \obv{converse_iff}.
        $c\mathrel{(\converse{\reflexiveTransitiveClosure{X}{R}}\circ
            \reflexiveTransitiveClosure{X}{R})}b$ by \auto{circ_iff}.
        Follows by \obv{joinability_relation}.
    \end{subproof}
    Follows by \auto{symmetric}.
\end{proof}

\begin{corollary}\label{joinable_reflexive}
    Let $R$ be a binary relation on $X$.
    Let $b\in X$.
    Then $b$ is joinable with $b$ in $X$ with respect to $R$.
\end{corollary}
\begin{proof}
    $\joinabilityRelation{X}{R}$ is reflexive on $X$
        by \obv{joinability_relation_is_reflexive}.
    Follows by \obv{joinable_on,reflexive_on}.
\end{proof}

\begin{corollary}\label{joinable_symmetric}
    Let $R$ be a binary relation on $X$.
    Let $b,c\in X$.
    Suppose $b$ is joinable with $c$ in $X$ with respect to $R$.
    Then $c$ is joinable with $b$ in $X$ with respect to $R$.
\end{corollary}
\begin{proof}
    $\joinabilityRelation{X}{R}$ is symmetric
        by \obv{joinability_relation_is_symmetric}.
    Follows by \auto{joinable_on,symmetric}.
\end{proof}

\begin{proposition}\label{confluent_at_implies_locally_confluent_at}
    Let $R$ be a binary relation on $X$.
    Let $a\in X$.
    Suppose $R$ is confluent at $a$ on $X$.
    Then $R$ is locally confluent at $a$ on $X$.
\end{proposition}
\begin{proof}
    $R\subseteq\reflexiveTransitiveClosure{X}{R}$
        by \obv{reflexive_transitive_closure_contains}.
    Follows by \obv{confluent_at_on,locally_confluent_at_on,subseteq}.
\end{proof}

\begin{proposition}\label{confluent_on_implies_locally_confluent_on}
    Let $R$ be a binary relation on $X$.
    Suppose $R$ is confluent on $X$.
    Then $R$ is locally confluent on $X$.
\end{proposition}
\begin{proof}
    Follows by \obv{confluent_on,locally_confluent_on,
        confluent_at_implies_locally_confluent_at}.
\end{proof}

\begin{proposition}\label{confluent_on_iff_peak_relation}
    Let $R$ be a binary relation on $X$.
    Then $R$ is confluent on $X$ iff
        $\reflexiveTransitiveClosure{X}{R}\circ
            \converse{\reflexiveTransitiveClosure{X}{R}}
        \subseteq\joinabilityRelation{X}{R}$.
\end{proposition}
\begin{proof}
    Let $T=\reflexiveTransitiveClosure{X}{R}$.
    Let $J=\joinabilityRelation{X}{R}$.
    $T$ is a binary relation on $X$ by notion.
    $\converse{T}$ is a binary relation on $X$ by notion.
    $T\circ\converse{T}$ is a binary relation on $X$ by notion.
    $T\circ\converse{T}$ is a relation by notion.
    $J$ is a binary relation on $X$ by notion.

    Show that if $R$ is confluent on $X$, then
        $T\circ\converse{T}\subseteq J$.
    \begin{subproof}
        Suppose $R$ is confluent on $X$.
        Show that for all $b,c$ such that
            $b\mathrel{(T\circ\converse{T})}c$ we have $b\mathrel{J}c$.
        \begin{subproof}
            Fix $b,c$.
            Suppose $b\mathrel{(T\circ\converse{T})}c$.
            Take $a$ such that $b\mathrel{\converse{T}}a$ and
                $a\mathrel{T}c$ by \obv{circ_iff}.
            $a\mathrel{T}b$ by \obv{converse_iff}.
            $T\subseteq X\times X$ by definition.
            $(a,b),(a,c)\in X\times X$ by \obv{subseteq}.
            Hence $a,b,c\in X$ by \obv{times_tuple_elim}.
            $R$ is confluent at $a$ on $X$ by \obv{confluent_on}.
            Thus $b$ is joinable with $c$ in $X$ with respect to $R$
                by \obv{confluent_at_on}.
            Follows by \obv{joinable_on}.
        \end{subproof}
        Show that for every $w\in T\circ\converse{T}$ we have $w\in J$.
        \begin{subproof}
            Fix $w\in T\circ\converse{T}$.
            Take $b,c$ such that $(b,c)=w$ by \auto{relation}.
            $b\mathrel{(T\circ\converse{T})}c$.
            Hence $b\mathrel{J}c$.
        \end{subproof}
        Follows by \obv{subseteq}.
    \end{subproof}

    Show that if $T\circ\converse{T}\subseteq J$, then
        $R$ is confluent on $X$.
    \begin{subproof}
        Suppose $T\circ\converse{T}\subseteq J$.
        Show that for every $a\in X$ we have $R$ is confluent at $a$ on $X$.
        \begin{subproof}
            Fix $a\in X$.
            Show that for all $b,c\in X$ such that
                $a\mathrel{T}b$ and $a\mathrel{T}c$ we have
                $b$ is joinable with $c$ in $X$ with respect to $R$.
            \begin{subproof}
                Fix $b,c\in X$.
                Suppose $a\mathrel{T}b$ and $a\mathrel{T}c$.
                $b\mathrel{\converse{T}}a$ by \obv{converse_iff}.
                Hence $b\mathrel{(T\circ\converse{T})}c$
                    by \auto{circ_iff}.
                Thus $b\mathrel{J}c$ by \obv{subseteq}.
                Follows by \obv{joinable_on}.
            \end{subproof}
            Follows by \obv{confluent_at_on}.
        \end{subproof}
        Follows by \obv{confluent_on}.
    \end{subproof}
    Follows by \auto.
\end{proof}

\begin{proposition}\label{locally_confluent_on_iff_peak_relation}
    Let $R$ be a binary relation on $X$.
    Then $R$ is locally confluent on $X$ iff
        $R\circ\converse{R}\subseteq\joinabilityRelation{X}{R}$.
\end{proposition}
\begin{proof}
    $R\circ\converse{R}$ is a binary relation on $X$ by notion.
    $R\circ\converse{R}$ is a relation by notion.
    $\joinabilityRelation{X}{R}$ is a binary relation on $X$ by notion.
    Show that if $R$ is locally confluent on $X$, then
        $R\circ\converse{R}\subseteq\joinabilityRelation{X}{R}$.
    \begin{subproof}
        Suppose $R$ is locally confluent on $X$.
        Show that for all $b,c$ such that
            $b\mathrel{(R\circ\converse{R})}c$ we have
            $b\mathrel{\joinabilityRelation{X}{R}}c$.
        \begin{subproof}
            Fix $b,c$.
            Suppose $b\mathrel{(R\circ\converse{R})}c$.
            Take $a$ such that $b\mathrel{\converse{R}}a$ and
                $a\mathrel{R}c$ by \obv{circ_iff}.
            $a\mathrel{R}b$ by \obv{converse_iff}.
            $R\subseteq X\times X$ by definition.
            $(a,b),(a,c)\in X\times X$ by \obv{subseteq}.
            Hence $a,b,c\in X$ by \obv{times_tuple_elim}.
            $R$ is locally confluent at $a$ on $X$
                by \obv{locally_confluent_on}.
            Thus $b$ is joinable with $c$ in $X$ with respect to $R$
                by \obv{locally_confluent_at_on}.
            Follows by \obv{joinable_on}.
        \end{subproof}
        Show that for every $w\in R\circ\converse{R}$ we have
            $w\in\joinabilityRelation{X}{R}$.
        \begin{subproof}
            Fix $w\in R\circ\converse{R}$.
            Take $b,c$ such that $(b,c)=w$ by \auto{relation}.
            $b\mathrel{(R\circ\converse{R})}c$.
            Hence $b\mathrel{\joinabilityRelation{X}{R}}c$.
        \end{subproof}
        Follows by \obv{subseteq}.
    \end{subproof}
    Show that if
        $R\circ\converse{R}\subseteq\joinabilityRelation{X}{R}$, then
        $R$ is locally confluent on $X$.
    \begin{subproof}
        Suppose $R\circ\converse{R}\subseteq\joinabilityRelation{X}{R}$.
        Show that for every $a\in X$ we have
            $R$ is locally confluent at $a$ on $X$.
        \begin{subproof}
            Fix $a\in X$.
            Show that for all $b,c\in X$ such that
                $a\mathrel{R}b$ and $a\mathrel{R}c$ we have
                $b$ is joinable with $c$ in $X$ with respect to $R$.
            \begin{subproof}
                Fix $b,c\in X$.
                Suppose $a\mathrel{R}b$ and $a\mathrel{R}c$.
                $b\mathrel{\converse{R}}a$ by \obv{converse_iff}.
                Hence $b\mathrel{(R\circ\converse{R})}c$
                    by \auto{circ_iff}.
                Thus $b\mathrel{\joinabilityRelation{X}{R}}c$
                    by \obv{subseteq}.
                Follows by \obv{joinable_on}.
            \end{subproof}
            Follows by \obv{locally_confluent_at_on}.
        \end{subproof}
        Follows by \obv{locally_confluent_on}.
    \end{subproof}
    Follows by \auto.
\end{proof}
